% ============================================================
% Book chapter: "Two Worlds: Models, Reality, and the
% Paradoxes of Confusing Them"
%
% This file is self-contained and compiles standalone with
% pdflatex. To fold it into a larger book manuscript, delete
% everything before \chapter{...} and after the final
% \end{document}, and \input{} or \include{} the remainder.
% ============================================================

\documentclass[11pt]{book}

\usepackage[utf8]{inputenc}
\usepackage[T1]{fontenc}
\usepackage{amsmath}
\usepackage{amssymb}
\usepackage[margin=1.15in]{geometry}
\usepackage{enumitem}

\title{Two Worlds}
\author{}
\date{}

\begin{document}

\chapter{Two Worlds: Models, Reality, and the Paradoxes of Confusing Them}
\label{chap:two-worlds}

\section{The Physicist's Two Worlds}

A physicist studying the real world does not study it directly for long. Within the first few pages of any serious inquiry, the raw material of experience --- falling apples, glowing filaments, deflected starlight --- is translated into symbols, and the symbols are set to work according to rules that owe nothing further to experience. What began as an investigation of the world in front of us becomes, almost immediately, an investigation of a second world: a world of variables, operators, and symmetries, built expressly to stand in for the first.

It is worth naming the two worlds plainly, because the rest of this chapter depends on keeping them apart. The \emph{real world} is where causality operates, where a measurement is taken and comes out one way rather than another, where an experiment either agrees with a prediction or it does not. The \emph{ideal world} is the world of the model: a formal system in which a claim is true or false by proof, not by observation, and in which the only kind of support a statement can have is logical support from other statements already accepted within the system.

This chapter is about a specific and recurring hazard in physics, and in the empirical sciences generally: the hazard of losing track, in the middle of an argument, of which of the two worlds one is currently reasoning in. Some notions belong properly to one world and cannot be carried across the border into the other without an explicit toll paid at customs --- a bridge, built and defended, tying some piece of the formalism to some piece of the world. When that toll goes unpaid, when a purely formal relation is treated as though it automatically carried empirical weight, or an empirical judgment is expected to fall directly out of formal symmetry, the result is paradox. Not a contradiction in nature, and not, usually, an error in the mathematics, but a mismatch between the two, smuggled through unexamined.

\section{Anatomy of the Two Worlds}

\subsection{The Ideal World}

The ideal world is a formal system: a set of primitives, axioms, and rules of inference, together with everything that can be derived from them. Consider Newton's second law together with an inverse-square force law. The claim that a body subject to such a force traces a conic section is not, at that point, something physicists go out and check case by case; it is a theorem. Whatever discovery is involved in celestial mechanics lies in locating a force law that fits the world --- the trajectory that follows from it is furnished by the ideal world for free, by proof rather than by observation.

The properties native to this world --- provability, deductive closure, logical equivalence, symmetry under a transformation group, consistency --- hold or fail to hold independently of whether the system describes anything real. A physicist can prove theorems about a Hamiltonian that corresponds to no laboratory system in this universe, and the proofs are no less valid for that. Two statements can be \emph{logically} equivalent, standing or falling together as a matter of syntax, while having utterly different standing with respect to the world the formalism is meant to describe.

\subsection{The Real World}

The real world traffics in brute occurrence. A measurement outcome is not entailed; it simply happens, once, at a place, at a time, as a matter of fact rather than proof. Causality, in the sense relevant here, is a real-world notion: it concerns what actually produces, constrains, or makes possible what else, not which propositions can be derived from which. Experimental support is, in the same sense, a fact about how the world's occurrences bear on our credence in a hypothesis --- not a fact about the hypothesis's internal logical architecture.

\subsection{The Bridge}

A theory earns no purchase on the world until its symbols are tied to procedures. The tradition running from Percy Bridgman's operationalism through Hans Reichenbach's \emph{coordinative definitions} to what are often called \emph{correspondence rules} is, at bottom, an attempt to describe exactly this tie.\footnote{See Reichenbach, \emph{Experience and Prediction} (1938), on coordinative definitions, and Carnap, \emph{Philosophical Foundations of Physics} (1966), on correspondence rules; Bridgman's operationalism is developed in \emph{The Logic of Modern Physics} (1927).} The symbol $m$ in $F = ma$ means nothing on its own; it becomes physics once a coordinative definition ties it to, say, the reading of a balance scale calibrated against a reference mass.

The bridge is where experimental support actually happens. It is not a real-world notion in the way causality is --- a bare fact about the world --- and it is not an ideal-world notion either --- a bare fact about the model. It is a relation between the two, and it is only ever as trustworthy as the bridge it depends on.

This is the master distinction of the chapter, and it lets us state its central claim precisely: most of what looks like a paradox of physics, or of the sciences more generally, is a fault in the bridge, mistaken for a fault in the world or in the mathematics. The remaining sections work through this claim in four settings: a purely logical case, in which the diagnosis is clearest (Hempel's paradox of confirmation); its mirror image (Goodman's new riddle of induction); two cases native to physics itself; and, finally, some remarks on what a disciplined use of the bridge would require.

\section{When the Seam Shows: Hempel's Paradox of Confirmation}

Take the generalization
\[
\text{``All ravens are black,''} \qquad R \to B,
\]
where $R$ denotes being a raven and $B$ denotes being black. In the ideal world this is logically equivalent, by contraposition, to
\[
\text{``All non-black things are non-ravens,''} \qquad \neg B \to \neg R.
\]
That equivalence is a fact about material implication. It costs nothing empirically; it holds in every model of the propositional calculus, ravens or no ravens, exactly as $2+2=4$ holds whether or not anyone is counting anything.

Now bring in a principle from the real-world side of the ledger, sometimes called \emph{Nicod's criterion}: an observed instance of ``$P$ and $Q$'' confirms the generalization ``all $P$ are $Q$.''\footnote{After Jean Nicod's early-twentieth-century work on induction and probability.} A black raven, on this criterion, confirms ``all ravens are black.'' But the criterion applies just as well to the contrapositive: a white shoe is a non-black non-raven, and so, by the same reasoning, confirms ``all non-black things are non-ravens.''

The paradox appears the moment these two observations are combined with an \emph{equivalence condition}: whatever confirms a hypothesis confirms anything logically equivalent to it. Since the two generalizations above are logically equivalent, the white shoe --- found in a wardrobe, nowhere near a raven --- ends up confirming ``all ravens are black.''

The diagnosis, in the vocabulary of this chapter, is exact. Logical equivalence is a syntactic relation, native to the ideal world, indifferent to the world's actual population of ravens, shoes, or anything else. Confirmation is not syntactic; it is a real-world-facing relation, sensitive to facts such as how many ravens there are compared with how many non-black objects there are. The equivalence condition quietly assumes that these two relations must move in lockstep --- that whatever the ideal world says about a pair of statements, the real-world evidential relation has no choice but to agree. That assumption is the mixing error, and it is what generates the paradox out of two premises that are, individually, difficult to deny.

Hempel's own response was to accept the conclusion and explain why it feels wrong.\footnote{C.\ G.\ Hempel, ``Studies in the Logic of Confirmation,'' \emph{Mind} 54 (1945).} The white shoe \emph{does} confirm ``all ravens are black,'' he held, but by an amount so small as to be practically unnoticeable, because non-black objects vastly outnumber ravens. This can be made precise with a Bayesian likelihood-ratio argument, developed independently by Janina Hosiasson-Lindenbaum and later I.\ J.\ Good.\footnote{J.\ Hosiasson-Lindenbaum, ``On Confirmation,'' \emph{Journal of Symbolic Logic} 5 (1940); I.\ J.\ Good, ``The White Shoe is a Red Herring,'' \emph{British Journal for the Philosophy of Science} 17 (1967).} For a hypothesis $H$ and evidence $E$, the shift in credence is governed by
\[
\frac{P(H \mid E)}{P(\neg H \mid E)} \;=\; \frac{P(E \mid H)}{P(E \mid \neg H)} \cdot \frac{P(H)}{P(\neg H)}.
\]
Sampling a random object and finding it to be a non-black non-raven raises the likelihood ratio only trivially, because such an object is nearly as likely to turn up whether or not every raven happens to be black; sampling a raven and finding it black raises the ratio far more sharply, because ravens are exactly the population the hypothesis constrains. The paradox dissolves once the equivalence condition is no longer asked to do real-world evidential work on its own, and the real-world base rates are allowed back into the calculation.

\section{The Mirror Image: Goodman's New Riddle of Induction}

Nelson Goodman's grue paradox runs the same diagnosis in the opposite direction.\footnote{N.\ Goodman, \emph{Fact, Fiction, and Forecast} (1955).} Define a predicate \emph{grue} to apply to an object if it is examined before some future date $t^\ast$ and found green, or examined at or after $t^\ast$ and found blue; define \emph{bleen} with the colors swapped. Every emerald examined so far is equally well described as green or as grue --- the observations to date cannot tell the two hypotheses apart. As far as the ideal world of formal predicates is concerned, green and grue are perfectly symmetric: nothing in first-order logic, or in the bare notion of a predicate, prefers one over the other.

Yet no physicist, asked to bet on the color of an emerald mined after $t^\ast$, treats the two hypotheses as equally good. What breaks the symmetry is a real-world, non-logical judgment about which predicates track a genuine causal regularity in the world and which do not --- what philosophers of science call \emph{projectibility}. Logic alone, the ideal world's native currency, cannot supply that judgment; it has to be imported from outside the formalism.

Where Hempel's paradox comes from smuggling ideal-world structure into a real-world relation (confirmation), Goodman's comes from expecting the ideal world's logic, unaided, to settle a question that only the real world's causal and nomic structure can decide. The two failures are mirror images of the same underlying mistake: forgetting that the bridge between the worlds has to be built, and cannot be assumed.

\section{Physics-Native Cases}

The two paradoxes above are drawn from logic and confirmation theory because the diagnosis is cleanest there. But physics supplies its own native cases of the same trap, and they are, if anything, higher stakes, because getting the diagnosis wrong can mean mistaking an artifact of the mathematics for a genuine feature of nature, or the reverse.

\subsection{Dirac's Negative-Energy Solutions}

The relativistic wave equation Dirac proposed in 1928 for the electron admits solutions of negative energy, on exactly the same formal footing as the positive-energy solutions that were the equation's original motivation.\footnote{P.\ A.\ M.\ Dirac, ``The Quantum Theory of the Electron,'' \emph{Proc.\ Roy.\ Soc.\ A} 117 (1928).} At first these were pure ideal-world artifacts: nothing in the experimental record of the late 1920s obviously corresponded to a particle of negative energy. The tempting mistake would have been to go one way or the other by fiat --- to declare the negative-energy branch physically meaningless because no bridge to it was yet in hand, or to declare it physically real by simple analogy with the positive-energy branch, with no bridge built at all.

Dirac did neither. He constructed an explicit bridge: the negative-energy states, assumed filled in the vacuum by the Pauli exclusion principle, would leave an observable hole when one was unoccupied, and that hole would behave as a particle of positive energy and positive charge.\footnote{P.\ A.\ M.\ Dirac, ``A Theory of Electrons and Protons,'' \emph{Proc.\ Roy.\ Soc.\ A} 126 (1930).} This is a coordinative definition in the fullest sense: a specific, checkable procedure by which a piece of the formalism is tied to a piece of the world. In 1932 Carl Anderson observed exactly such a particle in cloud-chamber tracks, and the positron passed from a mathematical artifact awaiting a bridge to a confirmed inhabitant of the real world.\footnote{C.\ D.\ Anderson, ``The Positive Electron,'' \emph{Phys.\ Rev.} 43 (1933).} The lesson is not that ideal-world features are always physical, nor that they never are; it is that the question can only be settled by building and testing the bridge, not by inspecting the formalism alone.

\subsection{Coordinate and Gauge Artifacts}

The opposite outcome is equally well attested. In Schwarzschild coordinates, the metric of a non-rotating black hole develops a singularity at the Schwarzschild radius $r = 2GM/c^2$. Read naively, this looks like a real-world catastrophe at a specific radius. It is not: switching to Eddington--Finkelstein or Kruskal--Szekeres coordinates removes the singularity entirely, revealing it as an artifact of the original coordinate chart. The one singularity with a real-world, coordinate-independent counterpart --- signaled by the divergence of curvature invariants --- sits at $r=0$, not at the horizon.

A related case is the appearance of phase or group velocities exceeding $c$ in regions of anomalous dispersion, or in certain evanescent-wave and tunneling setups. The formal quantity outruns the speed of light, but no signal, and no information, is thereby transmitted faster than $c$; the relevant real-world bound is on the front velocity of an actual signal, which relativity's causal structure still constrains. In both cases, an ideal-world feature --- a coordinate singularity, a superluminal phase velocity --- has no real-world causal counterpart at all, and treating it as though it did is Hempel's error transplanted into general relativity and optics.

\subsection{A Working Heuristic}

The contrast between these two cases suggests a heuristic, not a decision procedure: before granting an ideal-world feature real-world or causal status, exhibit the specific correspondence rule that licenses the move, in the manner Dirac did. Absent an exhibited bridge, the default is to treat the feature as a mathematical artifact of the chosen formalism --- a coordinate system, a choice of gauge, a branch of a solution set --- until a bridge is shown to exist. The heuristic cuts both ways: it also warns against dismissing an unfamiliar formal feature as ``merely mathematical'' before the possibility of a bridge has actually been checked.

\section{Toward a Discipline of the Bridge}

The thread running through every example in this chapter is the same: paradox tends to appear not inside either world, but at the unexamined seam between them. This is why the philosophical literature on scientific theories devotes so much attention to the nature of that seam. Reichenbach's coordinative definitions and Carnap's correspondence rules were early attempts to make the bridge an explicit part of a theory's structure rather than an unstated background assumption.\footnote{Reichenbach, \emph{Experience and Prediction} (1938); Carnap, \emph{Philosophical Foundations of Physics} (1966).} The later semantic, or model-theoretic, view of scientific theories, developed by Patrick Suppes and extended by Bas van Fraassen and others, goes further, treating a theory as a family of mathematical models together with a specification of which of their parts are meant to be compared with data --- building the bridge into the definition of what a theory even is.\footnote{P.\ Suppes, \emph{Representation and Invariance of Scientific Structures} (2002), collecting work from the 1960s onward; B.\ van Fraassen, \emph{The Scientific Image} (1980).}

A closely related but distinct problem, worth flagging so as not to conflate it with the one this chapter has been concerned with, is the Duhem--Quine problem: because a prediction typically follows only from a whole web of hypotheses acting together, a failed prediction --- a real-world disconfirmation --- cannot by itself identify which single hypothesis in the ideal-world deductive chain is at fault.\footnote{P.\ Duhem, \emph{The Aim and Structure of Physical Theory} (1906); W.\ V.\ O.\ Quine, ``Two Dogmas of Empiricism,'' \emph{Philosophical Review} 60 (1951).} That, too, is a problem that lives at the bridge, but it concerns the many-to-one relation between hypotheses and predictions rather than the mismatch between formal and evidential relations examined here; it is the natural subject of a chapter of its own.

For the working physicist, the moral of this chapter can be put as a short discipline, applied whenever an argument seems to move too easily from the page to the world or back again:

\begin{enumerate}[label=(\arabic*)]
    \item \textbf{Locate the claim.} Is it a claim about what follows from what within the formalism, or a claim about what the world's occurrences make more or less credible? These are different kinds of claim, and conflating them is the first step toward paradox.
    \item \textbf{Check the equivalence.} If two statements are equivalent under some formal operation --- contraposition, a change of coordinates, a choice of gauge --- ask whether the theory's actual correspondence rules preserve whatever real-world property is at stake, or only the formal one.
    \item \textbf{Demand the bridge.} Before treating an ideal-world feature as physically or causally real, or before dismissing one as ``merely formal,'' exhibit the specific rule that would settle the question, in the way Dirac's hole theory settled the question for negative-energy states.
\end{enumerate}

None of this makes the ideal world dispensable; quite the opposite. It is only by building an ideal world --- disciplined, deductively closed, answerable to proof --- that a physicist can say anything general at all about the real one. The paradoxes surveyed in this chapter are not an argument against building such worlds. They are a record of what happens when the crossing between them is made in silence, and a case for making it, every time, out loud.

\begin{thebibliography}{99}

\bibitem{anderson1933} Anderson, C.\ D. (1933). ``The Positive Electron.'' \emph{Physical Review}, 43(6), 491--494.

\bibitem{bridgman1927} Bridgman, P.\ W. (1927). \emph{The Logic of Modern Physics}. New York: Macmillan.

\bibitem{carnap1966} Carnap, R. (1966). \emph{Philosophical Foundations of Physics}. New York: Basic Books.

\bibitem{dirac1928} Dirac, P.\ A.\ M.\ (1928). ``The Quantum Theory of the Electron.'' \emph{Proceedings of the Royal Society A}, 117(778), 610--624.

\bibitem{dirac1930} Dirac, P.\ A.\ M.\ (1930). ``A Theory of Electrons and Protons.'' \emph{Proceedings of the Royal Society A}, 126(801), 360--365.

\bibitem{duhem1906} Duhem, P.\ (1906). \emph{The Aim and Structure of Physical Theory}. Paris: Chevalier et Rivi\`ere.

\bibitem{good1967} Good, I.\ J.\ (1967). ``The White Shoe is a Red Herring.'' \emph{British Journal for the Philosophy of Science}, 17(4), 322.

\bibitem{goodman1955} Goodman, N.\ (1955). \emph{Fact, Fiction, and Forecast}. Cambridge, MA: Harvard University Press.

\bibitem{hempel1945} Hempel, C.\ G.\ (1945). ``Studies in the Logic of Confirmation.'' \emph{Mind}, 54(213--214), 1--26, 97--121.

\bibitem{hosiasson1940} Hosiasson-Lindenbaum, J.\ (1940). ``On Confirmation.'' \emph{Journal of Symbolic Logic}, 5(4), 133--148.

\bibitem{quine1951} Quine, W.\ V.\ O.\ (1951). ``Two Dogmas of Empiricism.'' \emph{Philosophical Review}, 60(1), 20--43.

\bibitem{reichenbach1938} Reichenbach, H.\ (1938). \emph{Experience and Prediction}. Chicago: University of Chicago Press.

\bibitem{suppes2002} Suppes, P.\ (2002). \emph{Representation and Invariance of Scientific Structures}. Stanford: CSLI Publications.

\bibitem{vanfraassen1980} van Fraassen, B.\ C.\ (1980). \emph{The Scientific Image}. Oxford: Clarendon Press.

\end{thebibliography}

\end{document}
